% from ega0-20-fr.tex
\oldpage[0\textsubscript{IV-1}]{149}
$\operatorname{D\acute{e}r.\ cont}_A(C,L)\to
\operatorname{D\acute{e}r.\ cont}_A(B,L)$ est injectif ; il est par ailleurs
surjectif en vertu de \sref{0.20.7.2-fr}, donc il est bijectif ; il en résulte
que l'image de $v_{C/B/A}$ est nécessairement dense dans $\Omega^1_{C/A}$,
sans quoi le quotient de $\Omega^1_{C/A}$ par l'adhérence de
$\operatorname{Im}(v_{C/B/A})$ serait séparé et $\neq0$ et aurait donc un
quotient discret $L\neq0$, contrairement à ce qu'on vient de voir (compte tenu
de \sref{0.20.4.8-fr}). Par suite $v_{C/B/A}$, qui est un monomorphisme formel
en vertu de \sref{0.20.7.2-fr}, est aussi un épimorphisme formel
\sref{0.19.1.2-fr}, donc un bimorphisme formel.
\end{proof}

\begin{corollary}[20.7.7]
\label{0.20.7.7-fr}
Supposons que dans $B$ et dans $C$ le carré de tout idéal ouvert soit ouvert.
Si $C$ est une $B$-algèbre formellement étale, alors, pour tout idéal ouvert
$\mathfrak{R}_\lambda$ de $C$, l'homomorphisme \sref{0.20.7.3.1-fr} est
bijectif.
\end{corollary}

\begin{proof}
En effet, en vertu de \sref{0.19.1.1-fr}, cet homomorphisme est surjectif, et
il est injectif par \sref{0.20.7.3-fr}.
\end{proof}

\begin{proposition}[20.7.8]
\label{0.20.7.8-fr}
Soient $u:A\to B$ un homomorphisme continu d'anneaux topologiques,
$\mathfrak{R}$ un idéal de $B$, $C$ l'anneau topologique quotient
$B/\mathfrak{R}$, $v:B\to C$ l'homomorphisme canonique. Alors :
\begin{enumerate}
  \item[\rm(i)] Dans la suite exacte \sref{0.20.5.12.1-fr}
\[
  \mathfrak{R}/\mathfrak{R}^2\xrightarrow{\delta_{C/B/A}}
  \Omega^1_{B/A}\otimes_B C\xrightarrow{v_{C/B/A}}\Omega^1_{C/A}\to0
\]
l'homomorphisme $\delta_{C/B/A}$ est continu et l'homomorphisme $v_{C/B/A}$
est un morphisme strict de $C$-modules topologiques.
  \item[\rm(ii)] Pour que $\delta_{C/B/A}$ soit formellement inversible à
gauche \sref{0.19.1.5-fr}, il faut et il suffit que pour tout $C$-module
discret $L$ annulé par un idéal ouvert de $C$, l'homomorphisme canonique
\begin{equation}
\tag{20.7.8.1}
\label{0.20.7.8.1-fr}
  \operatorname{Exalcotop}_A(C,L)\to\operatorname{Exalcotop}_A(B,L)
\end{equation}
soit injectif.
\end{enumerate}
\end{proposition}

\begin{proof}
(i) La première assertion est évidente. D'autre part, l'homomorphisme
canonique $v\otimes v:B\otimes_A B\to C\otimes_A C$ est un morphisme strict
par définition de la topologie produit tensoriel, et on en déduit aussitôt
(cf.\ \sref{0.20.5.2-fr}) qu'il en est de même de $v_{C/B/A}$.

(ii) Dire que $\delta_{C/B/A}$ est formellement inversible à gauche signifie
que pour tout $C$-module discret $L$ annulé par un idéal ouvert de $C$,
l'homomorphisme canonique
\[
  \operatorname{Hom.\ cont}_C(\Omega^1_{B/A}\otimes_B C,L)\to
  \operatorname{Hom.\ cont}_C(\mathfrak{R}/\mathfrak{R}^2,L)
\]
est surjectif. Or, compte tenu de \sref{0.18.4.3-fr} et de
\sref{0.20.4.8-fr}, cela revient à dire que l'homomorphisme canonique
\[
  \operatorname{D\acute{e}r.\ cont}_A(B,L)\to
  \operatorname{Exalcotop}_B(C,L)
\]
est surjectif, et la conclusion résulte donc de la suite exacte
\sref{0.20.3.6.1-fr}.

Le fait que \sref{0.20.7.8.1-fr} soit injectif s'exprime encore de la manière
suivante, compte tenu de la définition des deux membres
\sref{0.18.4.1-fr} : pour tout idéal ouvert $\mathfrak{M}$ de $A$, tout idéal
ouvert $\mathfrak{N}$ de $B$ tel que $\mathfrak{M}B\subset\mathfrak{N}$, et
toute $(A/\mathfrak{M})$-extension $E$ de
$B/(\mathfrak{R}+\mathfrak{N})$ par un
$(B/(\mathfrak{R}+\mathfrak{N}))$-module $L$, telle que l'image réciproque de
$E$ par l'homomorphisme canonique
$B/\mathfrak{N}\to B/(\mathfrak{R}+\mathfrak{N})$ soit
$(A/\mathfrak{M})$-triviale, il existe un idéal ouvert
$\mathfrak{M}'\subset\mathfrak{M}$ de $A$,
